\documentclass{article}
\usepackage{amsmath, amssymb, amsthm}
\usepackage{hyperref}
\usepackage{geometry}
\geometry{margin=1in}

\title{Erd\H{o}s Problem \#1003}
\date{Last updated: 2025-09-08}
\author{Source: erdosproblems.com}

\begin{document}
\maketitle

\noindent\textbf{Status:} OPEN \quad \textbf{No prize}

\noindent\textbf{Tags:} number theory

\medskip

\noindent\textbf{Problem Statement:}

Are there infinitely many solutions to $\phi(n)=\phi(n+1)$, where $\phi$ is the Euler totient function?




Erd\H{o}s \cite{Er85e} says that, presumably, for every $k\geq 1$ the equation\[\phi(n)=\phi(n+1)=\cdots=\phi(n+k)\]has infinitely many solutions. Erd\H{o}s, Pomerance, and S\'{a}rk\"{o}zy \cite{EPS87} proved that the number of $n\leq x$ with $\phi(n)=\phi(n+1)$ is at most\[\frac{x}{\exp((\log x)^{1/3})}.\]See [946] for the analogous question with the divisor function, and [415] for a question about more general inequality patterns among consecutive values of $\phi$.




References

[EPS87] Erd\H{o}s, Paul and Pomerance, Carl and S\'{a}rk\"ozy, Andr\'{a}s, On locally repeated values of certain arithmetic functions. {III} . Proc. Amer. Math. Soc. (1987), 1--7.

[Er85e] Erd\H{o}s, P., Some problems and results in number theory . Number theory and combinatorics. Japan 1984 (Tokyo, Okayama and Kyoto, 1984) (1985), 65-87.

\noindent\textbf{Related OEIS sequences:} A001274


\bigskip
\noindent\small{Source: \url{https://www.erdosproblems.com/1003}}
\end{document}
